% Original PDF 31 / printed 21, complete chapter introduction.
% block: unbounded-operators-heading page:31 type:chapter
\clearpage\setcounter{section}{2}
\section{\BT{Basic notions of unbounded operators on Hilbert spaces}{Hilbert 空间上无界算子的基本概念}}\label{chap:unbounded-operators}

% block: unbounded-operators-introduction page:31 type:ordinary-paragraph
\BT{Much of the modern theory of partial differential equations and of their numerical discretization is most elegantly studied in the framework of Hilbert spaces and operators acting on them. We assume that the reader is familiar with the basic definitions and results concerning Hilbert spaces and bounded linear operators on them. Here we focus on closed unbounded linear operators with an emphasis on adjoints and duality. This is also a mature and well-developed theory, even if less studied. Most of the results here can be found in numerous references, such as [\BookNumber{23}, Chapter 2.6], [\BookNumber{56}, Chapter III, Section 5, and Chapter IV, Section 5.2], or [\BookNumber{71}, Chapter II, Section 6, and Chapter VII], often in the greater generality of Banach spaces or even more general topological vector spaces. However, for our applications, the Hilbert space context suffices and, moreover, it allows for a brief, but fairly self-contained, exposition. At the end of the chapter we return to the differential operators which arose in the last chapter as differentials of the de Rham complex and interpret them as closed unbounded operators on Hilbert space and compute their adjoints. A summary of the key points is given in Box~\ref{box:unbounded-operators-summary} at the end of this chapter.}{偏微分方程及其数值离散化的许多现代理论, 在 Hilbert 空间及其上算子的框架中研究最为简洁优美. 我们假定读者熟悉 Hilbert 空间及其上有界线性算子的基本定义和结论. 这里着重讨论闭的无界线性算子, 特别关注伴随与对偶. 这同样是一套成熟且发展完善的理论, 尽管受到的研究较少. 本章大多数结论可以在许多参考文献中找到, 如 [\BookNumber{23}, Chapter 2.6], [\BookNumber{56}, Chapter III, Section 5 及 Chapter IV, Section 5.2], 或 [\BookNumber{71}, Chapter II, Section 6 及 Chapter VII]; 这些文献常在更一般的 Banach 空间甚至更一般的拓扑向量空间中讨论. 不过, 对我们的应用而言, Hilbert 空间的背景已经足够, 而且还能使论述简短而较为自足. 在本章末尾, 我们将回到上一章中作为 de Rham 复形微分出现的微分算子, 将它们解释为 Hilbert 空间上的闭无界算子, 并计算其伴随. 主要内容概括于本章末尾的 Box~\ref{box:unbounded-operators-summary}.}
